\documentclass[tikz,border=10pt]{standalone}
\usepackage{ctex}
\usepackage{tikz}
\usetikzlibrary{shapes.geometric, arrows.meta, positioning, fit, backgrounds, calc}

% 定义颜色
\definecolor{phaseblue}{RGB}{227,242,253}
\definecolor{phaseborder}{RGB}{21,101,192}
\definecolor{milestoneyellow}{RGB}{255,253,231}
\definecolor{milestoneborder}{RGB}{249,168,37}
\definecolor{startgray}{RGB}{207,216,220}
\definecolor{startborder}{RGB}{69,90,100}
\definecolor{activitywhite}{RGB}{255,255,255}
\definecolor{activityborder}{RGB}{55,71,79}
\definecolor{feedbackorange}{RGB}{255,87,34}

% 定义样式
\tikzstyle{phase} = [
    rectangle, rounded corners=8pt,
    draw=phaseborder, line width=1.5pt,
    fill=phaseblue,
    minimum width=7cm, minimum height=5cm,
    font=\bfseries\small
]
\tikzstyle{activity} = [
    rectangle, rounded corners=4pt,
    draw=activityborder, line width=0.8pt,
    fill=activitywhite,
    text width=5.8cm, align=center,
    minimum height=1cm,
    font=\small
]
\tikzstyle{milestone} = [
    rectangle, rounded corners=6pt,
    draw=milestoneborder, line width=1.5pt,
    fill=milestoneyellow,
    text width=2.2cm, align=center,
    minimum height=2.2cm,
    font=\footnotesize
]
\tikzstyle{startend} = [
    rectangle, rounded corners=12pt,
    draw=startborder, line width=1.5pt,
    fill=startgray,
    text width=4.5cm, align=center,
    minimum height=1.2cm,
    font=\small
]
\tikzstyle{endbox} = [
    rectangle, rounded corners=12pt,
    draw=startborder, line width=1.5pt,
    fill=startgray,
    text width=4cm, align=center,
    minimum height=1.2cm,
    font=\small
]
\tikzstyle{arrow} = [
    -{Stealth[length=3mm]},
    line width=1pt,
    color=activityborder
]
\tikzstyle{feedback} = [
    -{Stealth[length=2.5mm]},
    line width=1pt,
    dashed,
    color=feedbackorange
]
\tikzstyle{phaselabel} = [
    font=\bfseries\footnotesize,
    color=phaseborder
]

\begin{document}
\begin{tikzpicture}[node distance=1cm]

% ===== 标题 =====
\node[font=\Large\bfseries, text width=20cm, align=center] (title) at (10, 0.8)
    {财务数据架构试点实施技术路线图};

% ===== 开始节点 =====
\node[startend] (start) at (0, -2) {项目启动\\接收财务架构与需求输入};

% ===== 第一阶段 =====
\node[phaselabel] (p1label) at (0, -4.5) {第一阶段：财务本体定义与架构映射 (2026.01)};
\node[activity] (a11) at (0, -5.8) {财务架构梳理与实体识别};
\node[activity] (a12) at (0, -7.4) {基于CIM的财务领域本体构建};
\node[activity] (a13) at (0, -9.0) {知识图谱关系映射定义};

\begin{scope}[on background layer]
    \node[phase, fit=(p1label)(a11)(a12)(a13), inner sep=12pt] (p1box) {};
\end{scope}

% ===== 里程碑1 =====
\node[milestone] (m1) at (5.5, -6.5) {里程碑 1月底\\财务本体库\\初步建模};

% ===== 第二阶段 =====
\node[phaselabel] (p2label) at (11, -4.5) {第二阶段：EAV存储构建与数据汇聚 (2026.02)};
\node[activity] (a21) at (11, -5.8) {财务场景EAV模型动态扩展设计};
\node[activity] (a22) at (11, -7.4) {流批一体数据接入适配};
\node[activity] (a23) at (11, -9.0) {多源异构财务数据自动化导入};

\begin{scope}[on background layer]
    \node[phase, fit=(p2label)(a21)(a22)(a23), inner sep=12pt] (p2box) {};
\end{scope}

% ===== 里程碑2 =====
\node[milestone] (m2) at (16.5, -6.5) {里程碑 2月底\\财务全量数据\\EAV化存储};

% ===== 第三阶段 =====
\node[phaselabel] (p3label) at (0, -13) {第三阶段：本体串联与语义对齐 (2026.03)};
\node[activity] (a31) at (0, -14.3) {EAV数据库与本体管理器串联集成};
\node[activity] (a32) at (0, -15.9) {基于混合打分算法的语义相似度计算};
\node[activity] (a33) at (0, -17.5) {财务-实物资产语义映射规则生成};

\begin{scope}[on background layer]
    \node[phase, fit=(p3label)(a31)(a32)(a33), inner sep=12pt] (p3box) {};
\end{scope}

% ===== 里程碑3 =====
\node[milestone] (m3) at (5.5, -15) {里程碑 3月底\\语义映射规则\\部署完成};

% ===== 第四阶段 =====
\node[phaselabel] (p4label) at (11, -13) {第四阶段：模型萃取与试点验证 (2026.04)};
\node[activity] (a41) at (11, -14.3) {面向多态业务的参数化切片技术研发};
\node[activity] (a42) at (11, -15.9) {财务审计/成本分析视点模型萃取};
\node[activity] (a43) at (11, -17.5) {试点场景应用与性能验证};

\begin{scope}[on background layer]
    \node[phase, fit=(p4label)(a41)(a42)(a43), inner sep=12pt] (p4box) {};
\end{scope}

% ===== 里程碑4 =====
\node[milestone] (m4) at (16.5, -15) {里程碑 4月底\\试点场景\\功能验收};

% ===== 结束节点 =====
\node[endbox] (end) at (16.5, -11) {项目目标达成\\建立统一数据与\\知识表达体系};

% ===== 连接线 =====
% 阶段内部
\draw[arrow] (a11) -- (a12);
\draw[arrow] (a12) -- (a13);
\draw[arrow] (a21) -- (a22);
\draw[arrow] (a22) -- (a23);
\draw[arrow] (a31) -- (a32);
\draw[arrow] (a32) -- (a33);
\draw[arrow] (a41) -- (a42);
\draw[arrow] (a42) -- (a43);

% 开始 -> 阶段1
\draw[arrow] (start) -- (p1box.north);

% 阶段1 -> 里程碑1 -> 阶段2
\draw[arrow] (p1box.east) -- (m1.west);
\draw[arrow] (m1.east) -- (p2box.west);

% 阶段2 -> 里程碑2
\draw[arrow] (p2box.east) -- (m2.west);

% 阶段1 -> 阶段3
\draw[arrow] (p1box.south) -- (p3box.north);

% 里程碑2 -> 结束
\draw[arrow] (m2.south) -- (end.north);

% 阶段3 -> 里程碑3 -> 阶段4
\draw[arrow] (p3box.east) -- (m3.west);
\draw[arrow] (m3.east) -- (p4box.west);

% 阶段4 -> 里程碑4 -> 结束
\draw[arrow] (p4box.east) -- (m4.west);
\draw[arrow] (m4.north) -- (end.south);

% ===== 反馈回路 (放在最外侧避免重叠) =====
\draw[feedback] (p4box.south) -- ++(0,-1.5) -- ++(-16,0) -- ++(0,12.5) -- (-5,-7.4) -- (a12.west)
    node[pos=0.15, below, font=\scriptsize, color=feedbackorange] {验证反馈与模型迭代优化};

\end{tikzpicture}
\end{document}
